% GENERATED FILE. Edit canonical lessons, then run npm run generate:books.
% canonical-source-hash: fnv1a64:23abcff94428f6be
% canonical-lessons: UR-C443-gham, UR-C443-dukh, UR-C443-mayus, UR-C443-kurhna, UR-C443-afsurda

\chapter{Sadness}
\label{ch:ur-sadness}
\hlchaptermodality{443}

\begin{chapteropening}
\textbf{By the end of this chapter:} \emph{I can say that someone was sad or let down.}

Complete the last lesson of chapter 443: I can say that someone was sad or let down.
\end{chapteropening}

\section[\texorpdfstring{\ur{غم} (gham)}{غم (gham)}]{\ur{غم} (gham) — grief}

\label{lesson:UR-C443-gham}



\begin{quote}
Before the new one: say the Urdu for a wonder, then the Urdu for an amazing thing.
\end{quote}

\subsection*{You'll want to know: \ur{غم}}
\textbf{\ur{غم}} — \emph{gham} — \textquotedblleft{}grief\textquotedblright{}.

\begin{quote}
\ur{اس} \ur{کے} \ur{دل} \ur{میں} \ur{بڑا} \ur{غم} \ur{ہے۔}
\end{quote}

\begin{quote}
\emph{us ke dil meṅ baṛā gham hai.}
\end{quote}

\begin{quote}
There is great grief in his heart.
\end{quote}

Masculine.

\begin{cousinweb}
\textbf{\ur{غم}} (\emph{gham}) is usually traced, through Persian, to Arabic \emph{ghamm}, \textquotedblleft{}grief, distress\textquotedblright{}.
\end{cousinweb}

\begin{grammarlens}[title={What you've built}]
One more word for describing what you lived and felt.
\end{grammarlens}

\subsection*{Your turn}
\begin{itemize}
\raggedright
  \item \emph{Say it:} \emph{gham}
  \item \emph{Say it:} \emph{gham}, once more
  \item \emph{Say it:} \emph{us ke dil meṅ baṛā gham hai.} — There is great grief in his heart.
  \item \emph{Recall:} say the Urdu for a wonder, then the Urdu for an amazing thing, then say \emph{gham} again
\end{itemize}

\begin{tcolorbox}[breakable,colback=teal!4,colframe=teal!35!black,title={Before you move on}]
What does \ur{غم} mean? (\textquotedblleft{}Grief\textquotedblright{}.) Say it once more.
\end{tcolorbox}
\section[\texorpdfstring{\ur{دکھ} (dukh)}{دکھ (dukh)}]{\ur{دکھ} (dukh) — sorrow; distress}

\label{lesson:UR-C443-dukh}



\begin{quote}
Before the new one: say the Urdu for an amazing thing, then the Urdu for grief.
\end{quote}

\subsection*{You'll want to know: \ur{دکھ}}
\textbf{\ur{دکھ}} — \emph{dukh} — \textquotedblleft{}sorrow; distress\textquotedblright{}.

\begin{quote}
\ur{دوست} \ur{کے} \ur{جانے} \ur{کا} \ur{مجھے} \ur{دکھ} \ur{ہے۔}
\end{quote}

\begin{quote}
\emph{dost ke jāne kā mujhe dukh hai.}
\end{quote}

\begin{quote}
I am sad that my friend has gone.
\end{quote}

Masculine.

\begin{cousinweb}
\textbf{\ur{دکھ}} (\emph{dukh}) is inherited from Sanskrit \emph{duḥkhá}, \textquotedblleft{}difficulty, pain\textquotedblright{} (CDIAL 6375).
\end{cousinweb}

\begin{grammarlens}[title={What you've built}]
One more word for describing what you lived and felt.
\end{grammarlens}

\subsection*{Your turn}
\begin{itemize}
\raggedright
  \item \emph{Say it:} \emph{dukh}
  \item \emph{Say it:} \emph{dukh}, once more
  \item \emph{Say it:} \emph{dost ke jāne kā mujhe dukh hai.} — I am sad that my friend has gone.
  \item \emph{Recall:} say the Urdu for an amazing thing, then the Urdu for grief, then say \emph{dukh} again
\end{itemize}

\begin{tcolorbox}[breakable,colback=teal!4,colframe=teal!35!black,title={Before you move on}]
What does \ur{دکھ} mean? (\textquotedblleft{}Sorrow; distress\textquotedblright{}.) Say it once more.
\end{tcolorbox}
\section[\texorpdfstring{\ur{مایوس} (māyūs)}{مایوس (māyūs)}]{\ur{مایوس} (māyūs) — disappointed; without hope}

\label{lesson:UR-C443-mayus}



\begin{quote}
Before the new one: say the Urdu for grief, then the Urdu for sorrow.
\end{quote}

\subsection*{You'll want to know: \ur{مایوس}}
\textbf{\ur{مایوس}} — \emph{māyūs} — \textquotedblleft{}disappointed; without hope\textquotedblright{}.

\begin{quote}
\ur{امتحان} \ur{کا} \ur{نتیجہ} \ur{دیکھ} \ur{کر} \ur{وہ} \ur{مایوس} \ur{ہوا۔}
\end{quote}

\begin{quote}
\emph{imtihān kā natījā dekh kar voh māyūs huā.}
\end{quote}

\begin{quote}
He was disappointed when he saw the exam result.
\end{quote}

\begin{cousinweb}
\textbf{\ur{مایوس}} (\emph{māyūs}) is usually traced, through Persian, to Arabic \emph{maʾyūs}, \textquotedblleft{}given up as hopeless\textquotedblright{}.
\end{cousinweb}

\begin{grammarlens}[title={What you've built}]
One more word for describing what you lived and felt.
\end{grammarlens}

\subsection*{Your turn}
\begin{itemize}
\raggedright
  \item \emph{Say it:} \emph{māyūs}
  \item \emph{Say it:} \emph{māyūs}, once more
  \item \emph{Say it:} \emph{imtihān kā natījā dekh kar voh māyūs huā.} — He was disappointed when he saw the exam result.
  \item \emph{Recall:} say the Urdu for grief, then the Urdu for sorrow, then say \emph{māyūs} again
\end{itemize}

\begin{tcolorbox}[breakable,colback=teal!4,colframe=teal!35!black,title={Before you move on}]
What does \ur{مایوس} mean? (\textquotedblleft{}Disappointed; without hope\textquotedblright{}.) Say it once more.
\end{tcolorbox}
\section[\texorpdfstring{\ur{کڑھنا} (kuṛhnā)}{کڑھنا (kuṛhnā)}]{\ur{کڑھنا} (kuṛhnā) — to fret inwardly; to grieve in silence}

\label{lesson:UR-C443-kurhna}



\begin{quote}
Before the new one: say the Urdu for sorrow, then the Urdu for disappointed.
\end{quote}

\subsection*{You'll want to know: \ur{کڑھنا}}
\textbf{\ur{کڑھنا}} — \emph{kuṛhnā} — \textquotedblleft{}to fret inwardly; to grieve in silence\textquotedblright{}.

\begin{quote}
\ur{وہ} \ur{اپنی} \ur{ناکامی} \ur{پر} \ur{کڑھتا} \ur{ہے۔}
\end{quote}

\begin{quote}
\emph{voh apnī nākāmī par kuṛhtā hai.}
\end{quote}

\begin{quote}
He frets over his failure.
\end{quote}

Past (he …): \ur{وہ} \ur{اس} \ur{بات} \ur{پر} \ur{کڑھا}, \emph{voh is bāt par kuṛhā}.

\begin{cousinweb}
\textbf{\ur{کڑھنا}} (\emph{kuṛhnā}) is inherited: CDIAL lists the Hindi cousin \emph{kuṛhnā}, \textquotedblleft{}to be annoyed, to grieve\textquotedblright{}, under Sanskrit \emph{kruddhá}, \textquotedblleft{}irritated\textquotedblright{} (CDIAL 3598).
\end{cousinweb}

\begin{grammarlens}[title={What you've built}]
One more word for describing what you lived and felt.
\end{grammarlens}

\subsection*{Your turn}
\begin{itemize}
\raggedright
  \item \emph{Say it:} \emph{kuṛhnā}
  \item \emph{Say it:} \emph{kuṛhnā}, once more
  \item \emph{Say it:} \emph{voh apnī nākāmī par kuṛhtā hai.} — He frets over his failure.
  \item \emph{Recall:} say the Urdu for sorrow, then the Urdu for disappointed, then say \emph{kuṛhnā} again
\end{itemize}

\begin{tcolorbox}[breakable,colback=teal!4,colframe=teal!35!black,title={Before you move on}]
What does \ur{کڑھنا} mean? (\textquotedblleft{}To fret inwardly; to grieve in silence\textquotedblright{}.) Say it once more.
\end{tcolorbox}
\section[\texorpdfstring{\ur{افسردہ} (afsurdā)}{افسردہ (afsurdā)}]{\ur{افسردہ} (afsurdā) — dejected; downcast}

\label{lesson:UR-C443-afsurda}



\begin{quote}
Before the new one: say the Urdu for disappointed, then the Urdu for to fret inwardly.
\end{quote}

\subsection*{You'll want to know: \ur{افسردہ}}
\textbf{\ur{افسردہ}} — \emph{afsurdā} — \textquotedblleft{}dejected; downcast\textquotedblright{}.

\begin{quote}
\ur{بارش} \ur{کے} \ur{دن} \ur{وہ} \ur{افسردہ} \ur{رہتی} \ur{ہے۔}
\end{quote}

\begin{quote}
\emph{bārish ke din voh afsurdā rahtī hai.}
\end{quote}

\begin{quote}
On rainy days she feels low.
\end{quote}

\begin{cousinweb}
\textbf{\ur{افسردہ}} (\emph{afsurdā}) is usually traced to Persian \emph{afsurda}, \textquotedblleft{}frozen, withered\textquotedblright{}: a downcast person is pictured as withered.
\end{cousinweb}

\begin{grammarlens}[title={What you've built}]
One more word for describing what you lived and felt.
\end{grammarlens}

\subsection*{Your turn}
\begin{itemize}
\raggedright
  \item \emph{Say it:} \emph{afsurdā}
  \item \emph{Say it:} \emph{afsurdā}, once more
  \item \emph{Say it:} \emph{bārish ke din voh afsurdā rahtī hai.} — On rainy days she feels low.
  \item \emph{Recall:} say the Urdu for disappointed, then the Urdu for to fret inwardly, then say \emph{afsurdā} again
\end{itemize}

\begin{tcolorbox}[breakable,colback=teal!4,colframe=teal!35!black,title={Before you move on}]
What does \ur{افسردہ} mean? (\textquotedblleft{}Dejected; downcast\textquotedblright{}.) Say it once more.
\end{tcolorbox}
